% !TEX root = ../main.tex
\section{Problema 6: Construcción de un Rectángulo}

\subsection{Descripción del Problema y Abstracción CSP}
% Descripción: colocación/empaquetamiento de n cuadrados de tamaños dados en un rectángulo de ancho W y alto H.
% Requerimientos: sin solapamiento, dentro de los límites del contenedor, todos los cuadrados deben ser utilizados.

\subsubsection{Parámetros y Datos de Entrada}
% n (número de cuadrados), tamaño de lado de cada cuadrado s_i, anchura W y altura H del rectángulo contenedor.

\subsubsection{Variables de Decisión y Dominios}
% Coordenadas de posición para cada cuadrado i: (x_i, y_i), típicamente la esquina inferior izquierda.
% Dominios: x_i in 0..(W - s_i), y_i in 0..(H - s_i).

\subsubsection{Restricciones del Modelo}
% Formalizar en lenguaje matemático:
% 1. Límites del contenedor (asegurados por los dominios o restricciones de cota).
% 2. No solapamiento entre pares de cuadrados:
%    Para cada par (i, j): (x_i + s_i <= x_j) \/ (x_j + s_j <= x_i) \/ (y_i + s_i <= y_j) \/ (y_j + s_j <= y_i)
%    O mediante predicados globales como diffn.

\subsection{Detalles de Implementación (\texttt{rectangulo.mzn})}

\subsubsection{Restricciones Redundantes y Rompimiento de Simetrías}
% Restricción redundante de área: sum_{i=1}^n (s_i^2) = W * H.
% Restricciones acumulativas 1D en X y en Y (cumulative).
% Rompimiento de simetrías geométricas: reflexiones horizontal/vertical y rotaciones (cuando W = H).
% Explicar qué simetrías se rompen y cómo recuperar soluciones simétricas.

\subsection{Estrategias de Búsqueda y Distribución}
% Anotaciones de búsqueda en MiniZinc adaptadas a problemas de empaquetamiento 2D.
% Heurísticas de ordenamiento (ej. ordenar cuadrados de mayor a menor tamaño).

\subsection{Pruebas y Análisis Experimental}
% Pruebas con la instancia del taller (cuadrados 3, 2, 2, 1, 1, 1 en rectángulo 5x4) y otras instancias.
% Estadísticas comparativas (nodos, fallos, tiempo).

\begin{table}[htbp]
    \centering
    \small
    \begin{tabular}{lcccc}
        \toprule
        Instancia & Configuración & Tiempo (s) & Nodos explorados & Fallos \\
        \midrule
        Taller ($5 \times 4$) & Base & -- & -- & -- \\
        Taller ($5 \times 4$) & Simetrías + Redundantes & -- & -- & -- \\
        Instancia 2 & Base & -- & -- & -- \\
        Instancia 2 & Optimizada & -- & -- & -- \\
        \bottomrule
    \end{tabular}
    \caption{Evaluación experimental de colocación de cuadrados.}
    \label{tab:rectangulo_resultados}
\end{table}

\subsection{Conclusiones del Ejercicio}
% Discusión sobre la efectividad del predicado global (diffn) frente a restricciones disyuntivas explícitas
% y el impacto del rompimiento de simetrías.
